\section{Literatür verisiyle yazılım uygulamaları}
\label{sec:spss-b09}

\textbf{Amaç ve veri.} \texttt{b09.csv} ile öğretim amaçlı
$H_0:\mu=10$ ve $H_1:\mu\ne10$ hipotezlerini sınayın
\citep{cortez2008data}. 10 puan, burada seçilen referans değeridir.
İki yönlü tek örneklem $t$ testinde anlamlılık düzeyi $\alpha=0{,}05$,
güven düzeyi yüzde 95 olarak belirlenmiştir.


\textbf{Ortak veri.} Doğrulanan B09 uygulamasında üç yazılım da
\nolinkurl{companion/spss/csv/b09.csv} dosyasını kullanmıştır. Dosya
395 satır ve 10 sütun içerir; kayıtların tamamı analiz edilir.
B06'nın seçim göstergeleri ve ağırlığı burada kullanılmaz.
Bu karşılaştırma Excel içe aktarmanın doğrulandığı anlamına gelmez.
Kodları \texttt{companion} klasörünü içeren paket kökünden çalıştırın;
veri sözlüğü ve hazırlık için bkz. sayfa~\pageref{sec:spss-guide}.

\subsection{IBM SPSS uygulaması}
\textbf{Başlamadan önce.} Aşağıdaki komut dosyası B09 CSV verisini
açar ve analiz eder. Menü yolunu kullanacaksanız önce veri açma ve
tanımlama komutlarını \texttt{EXECUTE} satırı dahil çalıştırın.
B08'den kalan tek satırlık özet yerine ham veriyi yeniden açın;
filtre, ağırlık ve dosya bölme kapalı olmalıdır. \texttt{GET DATA}
dosya açma hatası verirse önce \texttt{/FILE} yolunu düzeltin ve
kodu baştan çalıştırın. Veri açılmadan sonraki komutlar çalışmaz.

\begin{spssadimlar}
\spssadim{Önce \emph{Analyze $\to$ Descriptive Statistics $\to$ Explore} ile \texttt{g3} histogramını inceleyin. Referans değerini sonuçlara bakarak değiştirmeyin.}
\spssadim{\emph{Analyze $\to$ Compare Means $\to$ One-Sample T Test} yolunu açın. Yeni sürümlerde orta menü \emph{Compare Means and Proportions} olarak görünebilir.}
\spssadim{\texttt{g3} değişkenini \emph{Test Variable(s)} alanına taşıyın; \emph{Test Value} alanına 10 yazın.}
\spssadim{\emph{Options} altında güven düzeyini yüzde 95 bırakın; \emph{Continue $\to$ OK} seçin. Bu uygulamanın alternatif hipotezi iki yönlüdür.}
\spssadim{\emph{One-Sample Test} tablosunda \emph{t}, \emph{df} ve iki yönlü anlamlılık sütununu okuyun. Tek yönlü anlamlılık sütunu da varsa bu uygulamada onu kullanmayın.}
\end{spssadimlar}

{\raggedright
\noindent\textbf{Komutların görevi.} \texttt{TESTVAL=10} sıfır hipotezindeki ortalamayı, \texttt{VARIABLES=g3} test değişkenini belirtir. \texttt{CI(.95)} fark için yüzde 95 güven aralığı ister.\par}

\textbf{Uygulama kodu:} \texttt{b09}. Doğrulanan veri dosyası
\texttt{companion/spss} altındaki \texttt{csv/b09.csv} dosyasıdır.

\begin{lstlisting}[style=prismSPSS]
SET DECIMAL=DOT.
GET DATA /TYPE=TXT
 /FILE='companion/spss/csv/b09.csv'
 /ENCODING='UTF8'
 /ARRANGEMENT=DELIMITED /DELCASE=LINE
 /DELIMITERS="," /QUALIFIER='"'
 /FIRSTCASE=2
 /VARIABLES=id F8.0 school F8.0 sex F8.0 age F8.0
 studytime F8.0 absences F8.0 g1 F8.0 g2 F8.0
 g3 F8.0 pass10 F8.0.
FILTER OFF.
WEIGHT OFF.
SPLIT FILE OFF.
VARIABLE LABELS
 g3 'Yil sonu matematik notu'
 pass10 'Yil sonu notu 10 veya uzerinde'.
VALUE LABELS school 1 'GP' 2 'MS'.
VALUE LABELS pass10
 0 '10 puanin altinda' 1 '10 puan ve uzerinde'.
VARIABLE LEVEL
 id school sex pass10 (NOMINAL)
 studytime (ORDINAL)
 age absences g1 g2 g3 (SCALE).
FORMATS g3 (F14.9).
EXECUTE.
EXAMINE VARIABLES=g3
 /PLOT=HISTOGRAM /STATISTICS=DESCRIPTIVES
 /CINTERVAL=95.
T-TEST /TESTVAL=10 /VARIABLES=g3
 /CRITERIA=CI(.95).
\end{lstlisting}

\begin{outputguidebox}
\textbf{Paylaşılan SPSS çıktısıyla doğrulanan değerler.}
\emph{One-Sample Statistics} tablosunda $n=395$, ortalama \SPMatMean\ ve
standart hata \SPMatSE\ bulunmuştur. \emph{One-Sample Test} tablosunda
$t(394)=1{,}801$, çift yönlü $p=0{,}072$ gösterilmiştir. Yüzde 5 düzeyinde
$H_0$ reddedilmez; bu, ortalamanın kesinlikle 10 olduğunun veya farkın
önemsiz olduğunun kanıtı değildir. Test tablosundaki güven aralığı
\emph{ortalama farkı} içindir: Explore çıktısındaki ortalama aralığının
iki sınırından da 10 çıkarılır. Histogram, sıfır yığılmasını değerlendirmek
içindir; büyük örneklem bağımsızlık sorununu çözmez.
\end{outputguidebox}


\par\noindent\begin{minipage}{\linewidth}
\subsection{R uygulaması}
\label{sec:r-gercek-b09}

\texttt{mu=10} sıfır hipotezini belirler. \texttt{conf.int} ortalama için aralıktır; 10 çıkarılınca fark aralığı elde edilir.

\begin{lstlisting}[style=prismR]
veri <- read.csv("companion/spss/csv/b09.csv")
stopifnot(nrow(veri) == 395, ncol(veri) == 10,
          !anyNA(veri),
          all(veri$pass10 == as.integer(veri$g3 >= 10)))
sonuc <- t.test(veri$g3, mu=10,
               alternative="two.sided", conf.level=.95)
ozet <- c(
  n=nrow(veri), ort=mean(veri$g3), ss=sd(veri$g3),
  sh_ort=sd(veri$g3)/sqrt(nrow(veri)),
  t=unname(sonuc$statistic), df=unname(sonuc$parameter),
  p_iki_yonlu=sonuc$p.value, fark=mean(veri$g3)-10,
  ort_alt=sonuc$conf.int[1], ort_ust=sonuc$conf.int[2],
  fark_alt=sonuc$conf.int[1]-10,
  fark_ust=sonuc$conf.int[2]-10)
print(sonuc)
print(ozet, digits=15)
print(sum(veri$g3 == 0))
\end{lstlisting}

\textbf{Çıktıyı okuma.} Paylaşılan R çıktısı 395 satır, 10 sütun ve
her sütunda sıfır eksik değer gösterir. Test çıktısındaki
\texttt{p-value}, burada iki yönlü olasılıktır. Ayrıntılı özette
\texttt{ort\_alt} ve \texttt{ort\_ust} ortalama için;
\texttt{fark\_alt} ve \texttt{fark\_ust} ise ortalama eksi 10 için
sınırlardır. Sıfır notlu 38 kayıt analizde tutulmuştur.
\end{minipage}\par

\par\noindent\begin{minipage}{\linewidth}
\subsection{Python uygulaması}
\label{sec:python-b09}

\texttt{sonuc} iki yönlü testin istatistiğini ve $p$ değerini tutar.
\texttt{ddof=1} örneklem standart sapmasını seçer; ortalama aralığından
10 çıkarılarak SPSS fark aralığıyla karşılaştırılır.

\begin{lstlisting}[style=prismPythonApp]
import pandas as pd
import numpy as np
from scipy import stats

veri = pd.read_csv("companion/spss/csv/b09.csv")
assert veri.shape == (395, 10)
assert not veri.isna().any().any()
assert veri.pass10.eq((veri.g3 >= 10).astype(int)).all()
gozlem_sayisi = len(veri)
ortalama = veri.g3.mean()
standart_sapma = veri.g3.std(ddof=1)
standart_hata = standart_sapma / np.sqrt(gozlem_sayisi)
sonuc = stats.ttest_1samp(
    veri.g3, popmean=10, alternative="two-sided")
ortalama_alt, ortalama_ust = stats.t.interval(
    .95, df=gozlem_sayisi-1,
    loc=ortalama, scale=standart_hata)
ozet = pd.Series({
    "n": gozlem_sayisi, "ort": ortalama, "ss": standart_sapma,
    "sh_ort": standart_hata, "t": sonuc.statistic,
    "df": gozlem_sayisi-1, "p_iki_yonlu": sonuc.pvalue,
    "fark": ortalama-10,
    "ort_alt": ortalama_alt, "ort_ust": ortalama_ust,
    "fark_alt": ortalama_alt-10, "fark_ust": ortalama_ust-10
})
print(ozet.to_string(float_format=lambda deger: f"{deger:.13f}"))
print(int(veri.g3.eq(0).sum()))
\end{lstlisting}

\textbf{Çıktıyı okuma.} Paylaşılan Python çıktısında da 395 satır,
10 sütun, sıfır eksik değer ve 38 sıfır notlu kayıt vardır.
On iki özet değer, R sonuçları 13 ondalık basamağa yuvarlandığında
eşleşir. Her iki yazılımda ana veri korunur. Buradaki karşılaştırma
sayısal çıktılara aittir; R ve Python histogramları doğrulanmamıştır.
\end{minipage}\par

\subsection{Sonuçların karşılaştırılması ve ortak rapor}
12 Eylül 2026'da paylaşılan SPSS tabloları ve histogramı ile R ve
Python metin çıktıları incelenmiştir. Gözlem sayısı, ortalama,
standart sapma, standart hata, $t$, serbestlik derecesi, iki yönlü
$p$, ortalama farkı ve iki aralığın dört sınırından oluşan on iki
özet değer R ve Python'da 13 ondalık basamağa yuvarlandığında
eşleşmiştir. SPSS'in ilgili değerleri gösterilen yuvarlama
hassasiyetlerinde uyumludur. Proje CSV'sinden yapılan hesaplar
sonuçları destekler. Aşağıdaki tablolar çıktıların kitap düzenine
uyarlanmış özetidir; ekran görüntüsü değildir.

\begin{table}[H]
\centering
\caption{B09'da yıl sonu notunun doğrulanmış betimsel özeti}
\label{tab:b09-dogrulanmis-betimsel}
\begin{tabular}{@{}lr@{}}
\toprule
Ölçü & Değer \\
\midrule
Gözlem sayısı & 395 \\
Ortalama & 10,415190 \\
Standart sapma & 4,581443 \\
Ortalamanın standart hatası & 0,230517 \\
\bottomrule
\end{tabular}
\par\smallskip
{\footnotesize Ondalıklı değerler altı basamağa yuvarlanmıştır.
Standart sapma 394 böleniyle, standart hata $s/\sqrt{395}$ ile
hesaplanmıştır. \texttt{g3} için eksik değer yoktur; sıfır notlu
38 kayıt analizde korunmuştur.\par}
\end{table}

\begin{table}[H]
\centering
\caption{B09'da 10 puan referansına karşı iki yönlü tek örneklem testi}
\label{tab:b09-dogrulanmis-test}
\begin{tabular}{@{}lr@{}}
\toprule
Ölçü & Değer \\
\midrule
$t$ istatistiği & 1,801122 \\
Serbestlik derecesi & 394 \\
İki yönlü $p$ değeri & 0,072448 \\
Ortalama farkı ($\bar x-10$) & 0,415190 \\
Ortalamanın \%95 güven aralığı: alt & 9,961992 \\
Ortalamanın \%95 güven aralığı: üst & 10,868388 \\
Farkın \%95 güven aralığı: alt & $-0{,}038008$ \\
Farkın \%95 güven aralığı: üst & 0,868388 \\
\bottomrule
\end{tabular}
\par\smallskip
{\footnotesize Ondalıklı değerler altı basamağa yuvarlanmıştır;
ara hesaplarda yuvarlanmamış değerler kullanılmıştır. SPSS test
tablosu farkın aralığını, Explore tablosu ortalamanın aralığını
verir. Fark aralığının iki sınırına da 10 eklenince ortalama
aralığı elde edilir.\par}
\end{table}

\Needspace{15\baselineskip}
\begin{outputguidebox}
Tablo~\ref{tab:b09-dogrulanmis-test} için $p=0{,}072448>0{,}05$
olduğundan $H_0:\mu=10$ reddedilmez. Bu karar ortalamanın kesinlikle
10 olduğunu, iki değerin eşdeğerliğini veya farkın uygulamada
önemsizliğini kanıtlamaz. Ortalama aralığı 10'u, fark aralığı 0'ı
içerir; aynı iki yönlü test ve güven düzeyi için bu sonuçlar uyumludur.

Paylaşılan SPSS tablosunda \emph{One-Sided p} sütunu 0,036,
\emph{Two-Sided p} sütunu 0,072 gösterir. B09'un alternatif hipotezi
iki yönlü olduğundan ikinci sütun kullanılır. Yardım açıklamasındaki
tek yönlü anlamlılık vurgusu, bu uygulamanın kararını değiştirmez.
Sonuçlara bakarak testin yönünü değiştirmeyin; renk yerine hipotezi
ve ilgili sütunun başlığını esas alın.

SPSS histogramında sıfırda ayrı bir yığılma vardır; bu 38 kayıt
gerekçesiz biçimde silinmemiştir. Histogram bağımsızlığı veya temsil
gücünü doğrulamaz. Hesaplar öğretim amacıyla kullanılan bağımsız
gözlemler modeline dayanır; iki okuldan gelen verinin daha geniş bir
öğrenci evrenine genellenebilirliği, yazılımlar aynı sonucu verdiği
için kanıtlanmış olmaz.
\end{outputguidebox}

\Needspace{10\baselineskip}
\textbf{Raporlama örneği (doğrulanmış çıktılarla).}
``395 kayıtta yıl sonu notunun ortalaması 10,415190, standart sapması
4,581443 bulunmuştur. Öğretim amaçlı 10 puan referansına karşı iki
yönlü tek örneklem $t$ testinde yüzde 5 düzeyinde sıfır hipotezi
reddedilmemiştir ($t(394)=1{,}801122$, $p=0{,}072448$).
Ortalama farkı 0,415190 puan, farkın yüzde 95 güven aralığı
$[-0{,}038008;0{,}868388]$ bulunmuştur. Ortalamanın yüzde 95 güven
aralığı $[9{,}961992;10{,}868388]$ olarak hesaplanmıştır. Sıfır notlu
38 kayıt korunmuştur. Bağımsız gözlemler modeline dayanan bu sonuçlar,
iki okulla sınırlı veri kaynağının daha geniş bir öğrenci evrenine
genellenebilirliğini kanıtlamamaktadır.''


\textbf{Görev.} Farkın güven aralığına 10 ekleyerek ortalamanın
aralığını doğrulayın. Neden 0,036 yerine 0,072 sütununun kullanıldığını
hipotezlerle açıklayın. $H_0$'ın reddedilmemesi ile eşdeğerlik
kanıtı arasındaki ayrımı ve bağımsızlık varsayımının histogramdan
doğrulanamayacağını belirtin.